\documentclass[11pt]{article}
\usepackage[margin=25mm]{geometry}
\usepackage{amsmath,amssymb,amsthm}
\usepackage{booktabs}
\usepackage{hyperref}
\newtheorem{remark}{Remark}
\newtheorem{lemma}{Lemma}
\newtheorem{proposition}{Proposition}
\title{価値反復を確率的最短経路として捉えた優先順序伝播による高速化}
\author{vi\_reference 研究ノート}
\date{2026-06}
\begin{document}
\maketitle

\begin{abstract}
本家 ROS1 \texttt{value\_iteration} の u64 忠実移植を割引なし確率的最短経路（SSP）として定式化し、
固定点反復を Dijkstra 流の優先順序伝播へ置換した2ソルバ（\texttt{prio\_lc} 厳密／\texttt{prio\_ls} 近似）を
実装・計測した。優先順序は到達セルあたりの更新数を $O(N)$ 下界（$\approx 1$ 回）まで削減するが、
現行 CPU の wall-clock では per-update コストが増え、帯域最適な Jacobi 活性集合に及ばないことを示す。
\end{abstract}

\section{リファレンス実装の数式モデル}
状態 $s=(i_x,i_y,i_\theta)\in\mathcal{S}$、$|\mathcal{S}|=N=n_x n_y n_\theta$。索引
$\mathrm{idx}(s)=i_\theta+i_x n_\theta+i_y n_\theta n_x$。行動 $a\in\mathcal{A}$
（前進・後退・左右回転の6種）は、各 source $\theta$ について遷移分布
$\tau_a(i_\theta)=\{(\delta,p)\}$ を持つ。ここで $\delta=(d_{ix},d_{iy},d_{it})$、$d_{it}$ は
\emph{絶対}$\theta$ インデックス、$\sum_{(\delta,p)\in\tau_a(i_\theta)}p=B$、$B=2^{18}$。
$\tau_a$ はセルを $64^3$ にサブセル分割した決定論連続遷移
$T_a(x,y,\theta)=(x+\ell_a\cos\theta,\ y+\ell_a\sin\theta,\ \theta+\Delta_a)$ の集計である。

\paragraph{Bellman 作用素。}
\begin{equation}
(\mathcal{T}V)(s)=\min_{a\in\mathcal{A}}A(s,a),\quad
A(s,a)=\frac1B\sum_{(\delta,p)\in\tau_a(i_\theta)}p\bigl[V(s\oplus\delta)+g(s\oplus\delta)\bigr],
\label{eq:bellman}
\end{equation}
$s\oplus\delta=(i_x{+}d_{ix},\,i_y{+}d_{iy},\,d_{it}\bmod n_\theta)$。遷移先のいずれかが
範囲外または非 free なら $A(s,a)=\textsc{MaxCost}$。ステップコスト
$g(s')=\mathrm{penalty}(s')+\mathrm{local\_penalty}(s')\ge B$。固定点 $V^\star=\mathcal{T}V^\star$、
ゴール集合 $\mathcal{G}$（\texttt{final\_state}）は吸収 $V^\star|_{\mathcal{G}}=0$。
これは\textbf{割引なし確率的最短経路（SSP, $\gamma=1$）}であり、$g>0$・proper policy 存在
ゆえ固定点は到達可能セル上で一意。

\paragraph{実装の癖（remark）。}
\begin{remark}
式\eqref{eq:bellman}末尾の $\tfrac1B(\cdot)$ は u64 整数除算 $\gg18$ で切り捨てられる。
\end{remark}
\begin{remark}
$\sum p[\cdot]$ は u64 wrapping 加算。未到達セル $V=\textsc{MaxCost}$ 近傍は
オーバーフロー折返しで振動するが、到達可能セルの固定点には影響しない。
\end{remark}
\begin{remark}
本家は6種の sweep 順（行優先/列優先$\times$正逆＋分割2種）で Gauss--Seidel を行い、
直径律速を方向性で緩和する。$t_{\text{res}}=360/n_\theta$ は整数除算。
\texttt{State::from\_occupancy} の安全マージン penalty は線形 \texttt{pos} 境界のみ見る
行跨ぎバグを持つ。\texttt{valueFunctionWriter} は $\mathrm{total\_cost}/B$ の整数除算で出力。
\end{remark}

\section{計算量と高速化のレバー}
情報はゴールから1行動ステップ/スイープで伝播するため、素朴な Jacobi/Gauss--Seidel VI は
$O(\mathrm{diam}\cdot N)$ 仕事を要する（$\mathrm{diam}$＝行動ステップでのグラフ直径）。
活性集合（frontier）は波面のみ更新し $O(R\cdot N)$（$R$＝波面厚＝1セルが確定するまでの
平均更新回数）。一方、状態を値の昇順に\textbf{一度だけ}確定する優先順序伝播は
$O(N\log N)$（二分ヒープ）で、$\mathrm{diam}$ に依存しない。
\[
\text{Jacobi } O(\mathrm{diam}\cdot N)\ \gg\ \text{frontier } O(R\cdot N)\ \gtrsim\ \text{priority } O(N\log N).
\]
横長マップ（campus $14000\times800$, $\mathrm{diam}\approx2300$）でこの差が顕在化する。

\section{優先順序伝播ソルバ}
$V$ の tentative ラベルを $\ell(s)$（初期 $\infty=\textsc{MaxCost}$、$\mathcal{G}$ で $0$）とし、
最小ヒープで昇順に取り出す。cost-to-go では $V(s)$ が後続に依存するため、確定セル $s^\star$ は
それを outcome に持つ\emph{前駆}のラベルを改善し得る（逆$\theta$隣接 $\mathrm{rev}[i_t]$）。

\begin{verbatim}
priority_solve(label_setting):
  l <- MaxCost (G: 0); heap <- {(0,s) : s in G}
  while heap not empty:
    (lam, s*) <- pop-min
    if lam != l(s*): continue            # stale (遅延 decrease-key)
    settle 規則: A1=settle s* (確定済みなら skip) / A2=なし
    for 前駆 s in pred(s*), free, not in G (A1 は未確定のみ):
      l' <- (T l)(s)                     # 式(1)で再計算
      if l' < l(s): l(s) <- l'; push (l', s)
\end{verbatim}

\paragraph{A1: label-setting（\texttt{prio\_ls}、近似）.}
$s^\star$ を確定（settle）し二度と触れない。各到達セルを実質1回確定 $\Rightarrow$ 最小仕事量。
\paragraph{A2: label-correcting（\texttt{prio\_lc}、厳密）.}
settle せず、改善時は確定済みセルも再 relax。単調減少・下に有界ゆえ停止し、$V^\star$ に一致。
ゆえに本家と bit-exact。

\section{単調性とその破れ（A1 の近似誤差）}
\begin{proposition}
全行動の全 outcome が「確定済みかつ値 $\le V(s)$」なら、純粋な label-setting は厳密。
\end{proposition}
本問題ではこれが破れ得る：$A(s,a)$ は\emph{全 outcome の確定}を要するため、$s$ の最適行動が
「稀だが高値の outcome（後退・横滑り）＋多数の低値前進 outcome」を含むと、低値 outcome が
先に確定しても高値 outcome の確定が遅れ、その瞬間 $s$ のラベルが確定走査位置より\emph{後方}に
現れる（単調性違反）。
\begin{lemma}[違反の有界性]
$\tau_a$ はサブセル離散化由来で空間的に密集し、outcome の値の広がりは
$O(\max\|\delta\|\cdot \bar g)$ で有界。ゆえに違反量（後方挿入距離）はこの遷移スプレッドで
上から押さえられる。
\end{lemma}
よって A1 の誤差は局所的・小であり、決定論遷移（単一 outcome）では消える（実装テスト
\texttt{prio\_ls\_exact\_on\_deterministic\_transitions} で bit-exact を確認）。
本ノートでは違反量を \texttt{prio\_lc} の再処理回数 \texttt{repops} で実測する（\S5）。

\section{実験}
全実験は単スレッド・x86\_64 Linux (WSL2) で計測。コスト数式は全ソルバ共通（本家 u64 モデル）。

\subsection{house.pgm（$384\times384\times60$、実マップ）}
本家 ROS1 の固定点を ground truth とし、到達セルで比較。
\begin{center}
\begin{tabular}{lrrrrrl}
\toprule
ソルバ & elapsed[s] & 反復/pop & updates & 対本家 & RMSE & 方策一致 \\
\midrule
reference         & 23.49 & 61    & ---    & 1.08x  & 0.0000 & 100.00\% \\
frontier2d        & 7.82  & 64    & 33.4M  & 3.26x  & 0.0000 & 100.00\% \\
frontier2d\_pad   & 2.69  & 64    & 33.4M  & 9.46x  & 0.0000 & 100.00\% \\
frontier2d\_par   & 1.50  & 132   & 101.4M & 16.97x & 0.0000 & 100.00\% \\
\textbf{prio\_lc} & 21.89 & 2.88M & 3.24M  & 1.16x  & \textbf{0.0000} & \textbf{100.00\%} \\
\textbf{prio\_ls} & 11.41 & 2.84M & 3.16M  & 2.23x  & 187.49 & 99.58\% \\
\bottomrule
\end{tabular}
\end{center}
\texttt{prio\_lc} は本家と bit-exact（RMSE 0・方策 100\%）。更新数は frontier2d の 33.4M に対し
\textbf{3.24M（約 1/10）}だが、wall-clock は 21.9s と frontier2d の 7.8s より遅い（\S5.3）。
\texttt{prio\_ls} は値 RMSE 187 だが方策一致 99.58\%。

\subsection{直径レジーム（合成 free ストリップ、$W\times64\times60$）}
横長マップでゴールを左端に置き、直径 $\approx W/6$ を増大させる。
\begin{center}
\begin{tabular}{rlrrrrr}
\toprule
$W$ & ソルバ & elapsed[s] & pop/sweep & updates & upd/reach & repops \\
\midrule
512  & reference  & 25.16  & 77   & ---    & ---  & --- \\
512  & frontier2d & 2.75   & 96   & 7.97M  & 4.05 & --- \\
512  & prio\_ls   & 3.46   & 1.97M& 2.03M  & 1.03 & 0 \\
512  & prio\_lc   & 5.72   & 1.97M& 2.03M  & 1.03 & 1490 \\
\midrule
1024 & reference  & 97.52  & 127  & ---    & ---  & --- \\
1024 & frontier2d & 5.51   & 182  & 14.51M & 3.69 & --- \\
1024 & prio\_ls   & 7.03   & 3.93M& 4.12M  & 1.05 & 0 \\
1024 & prio\_lc   & 12.66  & 3.93M& 4.12M  & 1.05 & 1826 \\
\midrule
2048 & reference  & 340.86 & 215  & ---    & ---  & --- \\
2048 & frontier2d & 10.44  & 352  & 28.90M & 3.67 & --- \\
2048 & prio\_ls   & 14.23  & 7.86M& 8.33M  & 1.06 & 0 \\
2048 & prio\_lc   & 26.17  & 7.87M& 8.35M  & 1.06 & 4041 \\
\bottomrule
\end{tabular}
\end{center}
\texttt{prio\_ls} の近似精度（8x8 標準マップ、対 Reference 固定点）：empty RMSE 0.076 / 方策 99.87\%、
obstacle 2.71 / 99.68\%、sentinel 2.01 / 99.21\%。

\subsection{考察}
\paragraph{更新数の下界には到達する。}
優先順序伝播は \texttt{upd/reach} $\approx 1.03$--$1.06$ で、到達セルあたり\emph{ほぼ1回}の更新という
$O(N)$ 下界に到達する。活性集合 frontier の $\approx 3.7$--$4.05$ に対し約 $3.5\times$、素朴な
reference（house で 61 sweep）に対し桁違いに少ない。\textbf{数理的主張（優先順序 $\Rightarrow$ 最小更新数）は実証された。}

\paragraph{しかし wall-clock は追随しない（逆特性）。}
house で \texttt{prio\_lc} の更新数は frontier2d の $1/10$（3.24M vs 33.4M）だが、wall-clock は
$2.8\times$ 遅い（21.9 vs 7.8 s）。1 更新あたりコストは約 $29\times$ 高い。原因は、frontier の Jacobi
走査が逐次・キャッシュ整合（メモリ帯域律速）なのに対し、優先キューは二分ヒープの $O(\log N)$ pop と
前駆へのランダムアクセスでキャッシュミスを多発するため。\textbf{更新数で下界に達しても、現行ハードの
wall-clock では帯域最適 Jacobi が勝つ。}

\paragraph{直径律速は活性集合が既に解いている。}
reference は sweep 数が $77\to127\to215$ と直径に比例し、elapsed は $25\to97\to341$ s と
\emph{超線形}（$\propto \mathrm{diam}\times N$）。一方 frontier・優先順序はともに elapsed が $\propto N$
（幅倍化で約 $2\times$）。つまり「直径律速」は活性集合 frontier が既に解消しており、優先順序が
frontier に対して直径方向の追加利得を生むわけではない。優先順序の利得は\emph{定数倍の更新数削減}に
留まり、それは per-update コストで相殺される。

\paragraph{近似（\texttt{prio\_ls}）の精度。}
\texttt{prio\_ls} の値 RMSE はマップ複雑度とともに増大する（8x8 empty で 0.08、obstacle 2.71、
sentinel 2.01、house で 187）が、\emph{方策一致は常に高い}（99.2--99.9\%）。経路計画用途（方策）では
\texttt{prio\_ls} はほぼ厳密に等価、値の数値精度が要るなら \texttt{prio\_lc}（bit-exact）を使う。

\paragraph{単調性違反は極小 $\Rightarrow$ Dial 化が次の一手。}
\texttt{prio\_lc} の再処理 \texttt{repops} は pop 総数の $<0.1\%$（$1490/1.97\mathrm{M}\approx0.076\%$ 〜
$4041/7.87\mathrm{M}\approx0.051\%$）。\S4 の補題どおり違反は遷移スプレッドで有界であり、実測でも極小。
よってラベルが $\sim B$ 単位の整数である性質を使い、二分ヒープを\textbf{Dial（radix）バケット}に
置換すれば pop を $O(1)$ 化でき、per-update コストを下げて帯域最適 Jacobi との wall-clock 差を
縮められる見込みがある。これが本研究の指す具体的次手。

\section{結論}
本家 VI を割引なし SSP と捉え、固定点反復を Dijkstra 流の優先順序伝播へ置換した。厳密版
\texttt{prio\_lc} は本家と bit-exact、近似版 \texttt{prio\_ls} は方策をほぼ保ったまま更新数を $O(N)$
下界（到達セルあたり $\approx1$ 回）まで削減する。\textbf{「数式的にアルゴリズムの根本から」の
高速化レバー（最小更新数）は数理的に実在し実証された}が、現行 CPU の wall-clock では per-update
コスト（ヒープ＋ランダムアクセス、$\sim29\times$）が更新数削減を相殺し、既存の帯域最適 Jacobi
活性集合（\texttt{frontier2d\_par} で本家比 17x）に及ばない。直径律速は活性集合が既に解消済みで、
優先順序の追加利得は定数倍に留まる。ただし単調性違反が $<0.1\%$ と極小であることから、
\textbf{Dial radix バケット化}が wall-clock 改善の有望な次手として定量的に示された。
\end{document}
